\documentclass[11pt]{article}

\usepackage[margin=1in]{geometry}
\usepackage{graphicx}
\usepackage{amsmath,amssymb}
\usepackage{booktabs}
\usepackage{siunitx}
\usepackage{natbib}
\usepackage{microtype}
\usepackage{hyperref}

\title{An Atlas of Brain--Bone Sympathetic Neural Circuits Revealed by PRV152 Transneuronal Tracing}
\author{}
\date{}

\begin{document}
\maketitle

\begin{abstract}
The sympathetic nervous system (SNS) is an established regulator of bone metabolism, yet the identity of the central neurons that deliver SNS drive to bone is only partially understood in the mouse. Prior multi-synaptic tracing of rat femoral marrow outlined a limited hierarchical circuit, and extensive PRV tracing of adipose depots has defined central SNS outflow to fat, but a comprehensive, mouse femur-specific atlas has been missing. To fill this gap, we injected the attenuated, EGFP-expressing pseudorabies virus PRV152 (\SI{4.7e9}{pfu\per\milli\litre}; five 150\,nL loci) into the metaphysis and periosteum of the right femur of adult male mice ($N=6$). Following a 6-day survival interval, we processed every sixth 25\,$\mu$m brain section with anti-EGFP immunofluorescence and counted labeled neurons across the neuroaxis. Four of six animals produced uniform, analyzable labeling, while two were excluded for cloudy over-infection. Bone-surface application of PRV152 produced no EGFP signal in the paraventricular hypothalamic nucleus or raphe pallidus, and intermediolateral column labeling at T13--L2 confirmed the expected bone--SNS ganglia--IML--brain route. We identified 87 PRV152-positive nuclei, sub-nuclei and regions distributed across six brain divisions. The hypothalamus carried the largest number of labeled SNS neurons ($1177.25 \pm 62.75$), followed by midbrain and pons ($1065 \pm 22.39$), hindbrain medulla ($495.25 \pm 33.49$), forebrain ($237.5 \pm 15.08$), cerebral cortex ($104.75 \pm 4.64$), and thalamus ($65.25 \pm 7.78$). Dominant structures included the lateral hypothalamus, paraventricular hypothalamic nucleus and dorsomedial hypothalamus; periaqueductal gray and its lateral part, and the oral part of the pontine reticular nucleus; the raphe pallidus, raphe magnus and gigantocellular reticular nucleus; the medial preoptic nucleus, bed nucleus of the stria terminalis and ventral lateral septum; the primary somatosensory cortex hindlimb region together with the M1 and M2 motor cortex; and the periventricular fiber system and precommisural nucleus of the thalamus. The overlap of many femur-projecting hubs with previously described descending pain-modulatory nuclei identifies a candidate raphe magnus--periaqueductal gray axis for the central modulation of bone pain, and provides a neuroanatomical substrate for autonomic regulation of bone metabolism.
\end{abstract}

\section{Introduction}

Bone mass is maintained by balanced osteoblastic formation and osteoclastic resorption, and the sympathetic nervous system (SNS) is now firmly established as a modulator of that balance in the mouse. Genetic and pharmacological evidence has shown that \(\beta_2\)-adrenergic receptor signaling on osteoblasts suppresses bone formation and simultaneously promotes osteoclast differentiation through RANKL, while ablation of adrenergic signaling raises bone mass \citep{takeda2002leptin,elefteriou2005leptin}. Central leptin, acting on hypothalamic LepRb-expressing neurons, propagates an antiosteogenic signal to bone through this peripheral sympathetic relay \citep{ducy2000leptin,flak2016cns}. At the periphery, the skeleton is densely innervated by noradrenergic fibers immunoreactive for tyrosine hydroxylase, dopamine \(\beta\)-hydroxylase, or neuropeptide Y, and by a distinct population of vasoactive intestinal peptide-containing fibers that occupy the parenchyma of bone \citep{hohmann1986innervation,togari2010control}. Bone phenotypes in dopamine-transporter-deficient mice and norepinephrine uptake directly in the cortical compartment extend this picture to broader catecholaminergic tone \citep{bliziotes2000bone,zhu2018cortical}.

Despite these converging lines of evidence, the central architecture that delivers SNS drive specifically to bone remains incompletely mapped. The attenuated Bartha-derived strains of the swine alpha herpesvirus pseudorabies virus (PRV) provide the standard retrograde transneuronal tracer: PRV is endocytosed at peripheral axon terminals, replicates in post-synaptic somata, and then crosses synapses strictly in the retrograde direction, stepping serially through functionally connected neurons \citep{card2001transneuronal,ekstrand2008alpha,card1995influence}. Using PRV152, a Bartha-derived strain that expresses enhanced green fluorescent protein (EGFP) under the human cytomegalovirus immediate-early promoter, several groups have charted the central SNS outflow to white and brown adipose tissue, revealing a distributed network whose nodes include the medial preoptic region, paraventricular hypothalamus, lateral hypothalamus, periaqueductal gray, and medullary raphe, with partial depot-specific biases across the arcuate, dorsomedial, and suprachiasmatic nuclei \citep{bamshad1998central,bamshad1999cns,bartness1998innervation,bartness2014neural,ryu2014short,song2008mc4r,shrestha2010central,brito2007differential,adler2012neurochemical}. An earlier study in the rat outlined a hierarchical central circuit controlling SNS innervation of femoral marrow through lower thoracic and upper lumbar intermediolateral column (IML) preganglionic neurons and paravertebral chain ganglia. More recently, PRV152 and pharmacology together established that phosphodiesterase-5A (PDE5A)-containing neurons in discrete brain sites, including the locus coeruleus, raphe pallidus, and paraventricular hypothalamus, directly influence bone \citep{kim2020pde5a}.

Two questions have remained open. First, does the central SNS drive to bone simply recycle the circuitry used for adjacent peripheral targets such as adipose depots and bone marrow, or does the mouse femur recruit a partly private network? Second, do descending pain-modulatory structures, notably the periaqueductal gray (PAG) and the medullary raphe, deliver SNS outflow to bone, which would place the central control of bone metabolism and that of bone pain on at least partially overlapping anatomical substrates \citep{francois2017brainstem,heinricher2009descending,benarroch2012postural}? To address both questions, we built a comprehensive atlas of central SNS outflow to the mouse femur by injecting PRV152 into the metaphysis and periosteum, then performing serial anti-EGFP immunofluorescent analysis of the neuroaxis and cell counts in every sixth 25\,$\mu$m section.

Our contributions are fourfold. (i) We report the first comprehensive catalog of PRV152-labeled SNS neurons projecting from the central nervous system to the mouse femur, identifying 87 positive nuclei, sub-nuclei and regions across six brain divisions. (ii) We quantify per-division totals with means and standard errors, which rank the hypothalamus and the midbrain-and-pons as the dominant contributors and the cerebral cortex and thalamus as the least populated divisions. (iii) We identify the dominant hubs within each division, including the lateral hypothalamus, paraventricular hypothalamic nucleus and dorsomedial hypothalamus in the hypothalamus; the periaqueductal gray, lateral periaqueductal gray and oral part of the pontine reticular nucleus in the midbrain and pons; the raphe pallidus, raphe magnus and gigantocellular reticular nucleus in the medulla; the medial preoptic nucleus, bed nucleus of the stria terminalis and ventral lateral septum in the forebrain; the primary somatosensory cortex hindlimb region together with the M1 and M2 motor cortex in the neocortex; and the periventricular fiber system together with the precommisural nucleus in the thalamus. (iv) Combining these data with prior tracing in bone marrow and adipose tissue, we nominate a raphe magnus--periaqueductal gray axis as a candidate relay for descending control of bone pain. The remainder of the paper is organized as follows: Section~\ref{sec:methods} details the viral tracing, histology and counting; Section~\ref{sec:results} reports the quantitative atlas; and Section~\ref{sec:discussion} places these results in the context of prior SNS tracing and descending pain modulation.

\section{Methods}\label{sec:methods}

\subsection{Animals}
Adult male mice (approximately 3--4 months old) were single-housed on a 12\,h:12\,h light:dark cycle at $22\pm2$\,$^\circ$C with \emph{ad libitum} access to water and regular chow. All procedures were approved by the institutional animal care committee and performed in accordance with Public Health Service and United States Department of Agriculture guidelines.

\subsection{Viral construct and injections}
To identify central neurons projecting to bone, we used the attenuated Bartha-derived pseudorabies virus PRV152, which expresses EGFP under the human cytomegalovirus immediate-early promoter and spreads only in the retrograde direction across synapses after peripheral uptake \citep{card2001transneuronal,ekstrand2008alpha,card1995influence}. Six adult male mice ($N=6$) were anesthetized with isoflurane (2--3\% in oxygen) and the right femur--tibia joint was exposed. PRV152 (\SI{4.7e9}{pfu\per\milli\litre}) was injected at five loci (150\,nL per locus) distributed across the metaphysis and periosteum, which are known to be enriched in SNS innervation. At each locus the syringe was held in place for 60\,s to prevent efflux, and the incision was then closed with sterile sutures and wound clips. Topical nitrofurazone was applied to reduce infection risk. As a control for viral injection, PRV152 was also placed on the bone surface at the same titer; no EGFP signal was detected in pre-autonomic regions under this condition (Figure~\ref{fig:validation}). In addition, PRV152-labeled neurons in the intermediolateral cell column (IML) of the spinal cord at T13--L2 confirmed the expected bone--SNS ganglia--IML--brain route of infection \citep{bamshad1998central,ryu2014short}. All procedures involving PRV152 were conducted under Biosafety Level 2 containment.

\subsection{Histology and immunofluorescence}
Animals were sacrificed 6 days after the last PRV152 injection, a time point chosen on the basis of pilot studies that tracked viral progression to the brain. Mice were euthanized with carbon dioxide and perfused transcardially with 0.9\% heparinized saline, followed by 4\% paraformaldehyde in 0.1\,M phosphate-buffered saline (PBS; pH 7.4). Brains were post-fixed in the same fixative for 3--4\,h at 4\,$^\circ$C, transferred to 30\% sucrose in 0.1\,M PBS containing 0.1\% sodium azide, and stored at 4\,$^\circ$C until sectioning. Serial \SI{25}{\micro\metre} sections were cut on a freezing-stage sliding microtome and stored in 0.1\,M PBS with 0.1\% sodium azide until processing. For immunofluorescence, free-floating sections were rinsed in 0.1\,M PBS (2 $\times$ 15\,min), blocked for 30\,min in 10\% normal goat serum (NGS; Vector Laboratories) with 0.4\% Triton X-100 in 0.1\,M PBS, and incubated for 18\,h with a primary chicken anti-EGFP antibody (1:1000; Thermo Fisher Scientific A10262). Sections were then incubated for 2\,h at room temperature in an AlexaFluor-488-conjugated goat anti-chicken secondary antibody (1:700; Jackson Immunoresearch 103-545-155) with 2\% NGS and 0.4\% Triton X-100. Primary antibody omission and pre-adsorption with the immunizing peptide abolished immunoreactive staining in control sections. Sections were mounted on Superfrost Plus slides and cover-slipped with ProLong Gold Antifade (Thermo Fisher Scientific P36982).

\subsection{Imaging and quantitation}
Sections were imaged at $10\times$ and $20\times$ magnification on a Carl Zeiss Observer.Z1 fluorescence microscope with filters for AlexaFluor-488 and DAPI. Single-labeled PRV152 and DAPI images were overlaid in Zen and ImageJ. Cells positive for SNS PRV152 immunoreactivity were counted in every sixth brain section using the manual tag feature of Adobe Photoshop CS5.1, which eliminated double-counting of the same neuron. Counts were averaged across nucleus/sub-nucleus/region across all included animals. A standard mouse brain atlas (Paxinos and Franklin) was used to identify anatomical boundaries. Brightness, contrast and sharpness adjustments, removal of artifactual obstacles such as bubbles, and assembly into composite plates were performed in Photoshop for illustrative purposes only; no quantitative analyses were performed on adjusted images.

\subsection{Summary of injection parameters}
Table~\ref{tab:methods} summarizes the principal injection and tissue-processing parameters preserved verbatim from the experimental log.

\begin{table}[htbp]
\centering
\caption{Summary of injection, fixation and immunofluorescence parameters.}
\label{tab:methods}
\begin{tabular}{ll}
\toprule
Parameter & Value \\
\midrule
Animals & Adult male mice ($\sim$3--4 months), $N=6$ \\
Housing & Single, 12\,h:12\,h light:dark, $22\pm2$\,$^\circ$C \\
Viral tracer & PRV152 (Bartha-derived, EGFP+) \\
Titer & $4.7 \times 10^{9}$\,pfu/mL \\
Injection site & Right femur, metaphysis + periosteum \\
Number of loci & 5 \\
Volume per locus & 150\,nL \\
Syringe dwell time & 60\,s \\
Survival & 6 days post last injection \\
Perfusate 1 & 0.9\% heparinized saline \\
Perfusate 2 & 4\% paraformaldehyde in 0.1\,M PBS (pH 7.4) \\
Post-fix & 3--4\,h at 4\,$^\circ$C \\
Cryoprotection & 30\% sucrose in 0.1\,M PBS (0.1\% sodium azide) \\
Section thickness & \SI{25}{\micro\metre} \\
Blocking & 10\% NGS, 0.4\% Triton X-100, 30\,min \\
Primary antibody & chicken anti-EGFP, 1:1000 (Thermo Fisher A10262), 18\,h \\
Secondary antibody & AlexaFluor-488 goat anti-chicken, 1:700 (Jackson 103-545-155), 2\,h \\
Mounting medium & ProLong Gold Antifade (Thermo Fisher P36982) \\
Imaging & Zeiss Observer.Z1, 10$\times$ and 20$\times$, AF-488 and DAPI filters \\
Counting rule & every sixth brain section, manual Photoshop tags \\
\bottomrule
\end{tabular}
\end{table}

\section{Results}\label{sec:results}

\subsection{Validation of PRV152 transneuronal tracing from femoral bone}

We began by verifying that PRV152 delivered to bone propagates to the central nervous system through the expected SNS route rather than through non-specific diffusion. After inoculation, animals remained asymptomatic until day 5 post-inoculation, after which subjects showed occasional loss of body weight, decreased mobility, or an ungroomed coat; we euthanized each mouse for histology as symptoms emerged. Of the six injected mice, four exhibited uniform PRV152 infection across the neuroaxis from the hindbrain to the forebrain and were included in the quantitative atlas. Two animals showed cloudy plaque over-infection consistent with viral saturation and were excluded from downstream analyses. Consistent with our prior tracing of adipose depots, unilateral right-femur PRV152 injection produced bilateral brain labeling with no discernible ipsilateral dominance \citep{bamshad1998central,bartness1998innervation,ryu2014short}.

To confirm that central labeling depended on SNS uptake at the injected bone compartment and not on surface diffusion, we carried out two complementary controls (Figure~\ref{fig:validation}). First, when PRV152 at the same titer was placed on the bone surface rather than injected into the periosteum or metaphysis, no EGFP signal was detected in the paraventricular hypothalamic nucleus (PVH), which is known to harbor sympathetic pre-autonomic neurons, nor in the medullary raphe pallidus (RPa). Second, PRV152 injected directly into the periosteum or metaphysis produced robust EGFP immunoreactivity in the PVH and, critically, labeled neurons in the intermediolateral cell column (IML) of the spinal cord at T13--L2, the preganglionic segments consistent with a bone--SNS ganglia--IML--brain route of infection. Representative photomicrographs in the source dataset additionally illustrate PRV152 labeling in the periaqueductal gray (PAG) of the midbrain and pons, the RPa of the medulla, the lateral hypothalamus (LH), the medial preoptic nucleus medial part (MPOM) of the forebrain, the primary somatosensory cortex hindlimb region (S1HL) of the cerebral cortex, and the periventricular fiber system (pv) of the thalamus. Together these controls validate that the transneuronal spread of PRV152 from the femur reports bona fide SNS connectivity and not artefactual surface uptake. Having established that PRV152 from the femur traces a specific retrograde SNS pathway, we next quantified how its labeled neurons distribute across the neuroaxis.

\subsection{Distribution of PRV152-labeled neurons across six brain divisions}

Across the four analyzable animals, we identified 87 PRV152-positive nuclei, sub-nuclei and regions that collectively partitioned into six classical brain divisions: hypothalamus, midbrain and pons, hindbrain medulla, forebrain, cerebral cortex, and thalamus. To quantify the relative contribution of each division to bone-directed SNS outflow, we summed the PRV152-labeled neurons within every nucleus assigned to a given division and averaged across animals (Figure~\ref{fig:counts}, Table~\ref{tab:counts}). The hypothalamus contained the largest population of bone-labeled SNS neurons ($1177.25 \pm 62.75$), followed in descending order by the midbrain and pons ($1065 \pm 22.39$), the hindbrain medulla ($495.25 \pm 33.49$), the forebrain ($237.5 \pm 15.08$), the cerebral cortex ($104.75 \pm 4.64$), and the thalamus ($65.25 \pm 7.78$).

The distribution across divisions is monotonic: the hypothalamus and the combined midbrain and pons each harbor roughly an order of magnitude more labeled SNS neurons than the thalamus, and together account for more than 70\% of all bone-directed PRV152-positive neurons. The hindbrain medulla occupies an intermediate tier, consistent with its position in the canonical pre-autonomic cascade, while the cerebral cortex and thalamus each contribute a small but non-zero population of bone-projecting neurons. These totals establish a quantitative anatomical hierarchy: the primary source of femur-directed SNS outflow is the hypothalamus--midbrain/pons corridor, with meaningful but secondary contributions from the medulla and a modest cortical-thalamic footprint.

\begin{table}[htbp]
\centering
\caption{PRV152-labeled SNS neurons projecting to the mouse femur, summed within each brain division (mean $\pm$ SEM; $N=4$ analyzable animals).}
\label{tab:counts}
\begin{tabular}{lcc}
\toprule
Brain division & Mean count & SEM \\
\midrule
Hypothalamus & 1177.25 & 62.75 \\
Midbrain and pons & 1065 & 22.39 \\
Hindbrain medulla & 495.25 & 33.49 \\
Forebrain & 237.5 & 15.08 \\
Cerebral cortex & 104.75 & 4.64 \\
Thalamus & 65.25 & 7.78 \\
\bottomrule
\end{tabular}
\end{table}

\subsection{Dominant PRV152-positive nuclei within each division}

The aggregate counts above conceal important structure at the level of individual nuclei. We therefore identified, within each division, the nuclei with both the highest percentages and the highest absolute counts of PRV152-labeled neurons (Table~\ref{tab:hotspots}, Figure~\ref{fig:subnuclei}).

Within the hypothalamus, PRV152-positive neurons were distributed across 25 nuclei, sub-nuclei and regions. The lateral hypothalamus (LH), the paraventricular hypothalamic nucleus (PVH) and the dorsomedial hypothalamus (DM) were the dominant hubs by both percentage and absolute count. The LH and PVH in particular ranked among the most heavily labeled structures in the entire atlas. The midbrain and pons contributed 18 positive nuclei; here the periaqueductal gray (PAG), its lateral component (LPAG) and the oral part of the pontine reticular nucleus (PnO) were the most intensely labeled. The hindbrain medulla contributed 23 positive nuclei, with the raphe pallidus (RPa), raphe magnus (RMg) and gigantocellular reticular nucleus (Gi) carrying the largest percentages and counts of PRV152-labeled neurons. In the forebrain, 15 positive nuclei were identified, of which the medial preoptic nucleus, medial part (MPOM), the bed nucleus of the stria terminalis (BST) and the ventral part of the lateral septal nucleus (LSV) were the most heavily labeled. The cerebral cortex contributed only three positive regions, namely the primary somatosensory cortex hindlimb region (S1HL), the secondary motor cortex (M2) and the primary motor cortex (M1); of these, S1HL and M2 carried both the highest percentages and absolute counts. Finally, three thalamic sites were labeled, with the periventricular fiber system (pv) and the precommisural nucleus (PrC) carrying the highest counts.

\begin{table}[htbp]
\centering
\caption{Number of PRV152-positive structures and dominant hubs per brain division. ``Count of positive structures'' reflects nuclei, sub-nuclei and regions jointly.}
\label{tab:hotspots}
\begin{tabular}{lcl}
\toprule
Division & Count of positive structures & Top PRV152-positive hubs \\
\midrule
Hypothalamus & 25 & LH, PVH, DM \\
Midbrain and pons & 18 & PAG, LPAG, PnO \\
Hindbrain medulla & 23 & RPa, RMg, Gi \\
Forebrain & 15 & MPOM, BST, LSV \\
Cerebral cortex & 3 & S1HL, M2, M1 \\
Thalamus & 3 & pv, PrC \\
\bottomrule
\end{tabular}
\end{table}

The anatomy implied by these hot-spots is coherent. The most heavily labeled hypothalamic nuclei (LH, PVH, DM) are pre-autonomic centers whose roles in feeding, neuroendocrine control and autonomic integration are well established, and which express leptin receptors \citep{flak2016cns}. The dominant midbrain and medullary hubs (PAG, RMg, RPa) are the canonical descending pain-modulatory system \citep{francois2017brainstem,heinricher2009descending}. The forebrain MPOM and BST, together with the S1HL-centered cortical contribution, introduce a limbic and sensorimotor layer that plausibly integrates interoceptive and nociceptive signals from bone with autonomic outflow. The relatively small thalamic contribution, confined to pv and PrC, indicates that the femur is not a major target of classical thalamic SNS drive but is nonetheless sampled by midline/pre-commissural thalamic relays.

\subsection{A candidate brain--bone pain-relevant circuit}

Across the three results subsections, the PRV152 dataset recapitulates two observations with important interpretive consequences. First, bone-directed SNS outflow is largely \emph{distributed}: 87 structures across six divisions participate, so no single node is necessary or sufficient. Second, the strongest hot-spots are concentrated in structures that are \emph{simultaneously} autonomic and pain-modulatory. The PAG and the medullary raphe (RMg, RPa) together with the PVH form an anatomically coherent candidate axis (Figure~\ref{fig:circuit}): PVH pre-autonomic neurons project to the PAG; PAG, in turn, interfaces with raphe pallidus and raphe magnus; raphe projections terminate in the dorsal horn of spinal gray matter, where they regulate enkephalin-mediated gating of substance P-driven nociceptive transmission \citep{francois2017brainstem,heinricher2009descending,benarroch2012postural}; and preganglionic SNS output is then delivered to the femur via IML neurons at T13--L2 and paravertebral chain ganglia. This arrangement is consistent with, and extends, the hierarchical femoral marrow circuitry and previous PDE5A-bone mapping \citep{kim2020pde5a}.

\section{Discussion}\label{sec:discussion}

We present a comprehensive atlas of the central SNS circuitry that innervates the mouse femur, built from retrograde PRV152 tracing and manual counting across the neuroaxis. Across four analyzable animals, 87 PRV152-positive structures spanning six brain divisions contributed to bone-directed SNS outflow, with the hypothalamus ($1177.25 \pm 62.75$) and the midbrain-and-pons ($1065 \pm 22.39$) dominating the distribution. The dominant hubs within each division---LH, PVH, DM in the hypothalamus; PAG, LPAG, PnO in the midbrain and pons; RPa, RMg, Gi in the medulla; MPOM, BST, LSV in the forebrain; S1HL, M2, M1 in the cerebral cortex; and pv, PrC in the thalamus---together describe a distributed network that places bone-directed SNS neurons within structures classically associated with autonomic integration, descending pain modulation, and limbic-sensorimotor processing. These findings establish, for the first time in the mouse femur, the kind of detailed central map that has long existed for adipose depots and, in abbreviated form, for rat femoral marrow.

The anatomy we report is consistent with, yet substantially extends, prior work on bone and on adjacent peripheral targets. SNS outflow to white and brown adipose depots, traced by multiple groups with PRV152 and related reagents, converges on MPOM, PVH, LH, PAG, the oral part of the pontine reticular nucleus, and nuclei of the medullary raphe \citep{bamshad1998central,bamshad1999cns,bartness1998innervation,bartness2014neural,ryu2014short,song2008mc4r,shrestha2010central,brito2007differential,adler2012neurochemical}. The degree of overlap between these nodes and our femur atlas is striking: many of the same structures emerge, although the relative weight across divisions differs. A parsimonious reading is that bone is not innervated by a uniquely private neuronal population but by \emph{overlapping} SNS circuitry shared with other peripheral targets, with site-specific variation in the strength with which each central hub engages each target. This overlap explains why manipulations of generic SNS signaling---ranging from \(\beta_2\)-adrenergic receptor ablation \citep{takeda2002leptin,elefteriou2005leptin} to central leptin infusion \citep{ducy2000leptin,flak2016cns} to central PDE5A inhibition \citep{kim2020pde5a}---can simultaneously affect bone mass and other metabolic variables.

The second conceptual implication of our atlas is that femur-directed SNS neurons cluster heavily in canonical descending pain-modulatory structures. The PAG receives afferents from spinal cord, parabrachial nucleus and medullary raphe, and is a critical node for descending modulation of nociceptive transmission; the raphe magnus, in particular, has a well-documented role in the central modulation of noxious stimuli \citep{francois2017brainstem,heinricher2009descending,benarroch2012postural}. Projections from the medullary raphe, including RPa and RMg, terminate in the dorsal horn of spinal gray matter, where they regulate enkephalinergic gating of substance P-mediated nociception \citep{francois2017brainstem}. Our observation that PRV152 heavily labels PAG/LPAG, RMg and RPa after femoral injection therefore places bone-directed SNS neurons within the same anatomical scaffold that sets the gain on nociceptive transmission. The parsimonious interpretation is a candidate RMg--PAG ascending axis (relayed via PVH) that both delivers sympathetic drive to bone and modulates bone pain. This aligns with clinical observations that PAG-centered circuits participate in affective and autonomic control of pain, and offers a handle on bone pain that is tractable both pharmacologically and via neuromodulation.

Our data also bear on the leptin--hypothalamus--bone axis. Two of the three dominant hypothalamic hubs we identify (LH and PVH) are canonical leptin-responsive regions, and the third (DM) also expresses leptin receptors \citep{flak2016cns}. The LH dominates femur-directed outflow in our dataset, whereas prior rat marrow work emphasized the PVH. A reasonable interpretation is that PVH provides the principal marrow-directed SNS drive while LH and DM couple the anti-osteogenic actions of leptin more directly to the femur. This is consistent with classical demonstrations that leptin suppresses bone formation through a hypothalamic relay \citep{ducy2000leptin,takeda2002leptin}. A similar argument applies to the dopaminergic system: dopamine-transporter-deficient mice are osteopenic \citep{bliziotes2000bone}, and catecholaminergic tone buffered at the cortical compartment \citep{zhu2018cortical} would be influenced by descending inputs from the hubs we identify. Our atlas thus provides an anatomical substrate on which multiple neuroendocrine and neurochemical modulators of bone converge.

Several limitations should be acknowledged. First, our sample is intentionally small---six animals injected, four analyzable---because PRV152 tracing requires individually inspected, non-over-infected specimens and is not amenable to the sample sizes typical of behavioral or molecular work. Second, all mice were adult males; sex-specific differences in the projection pattern are plausible given evidence for sexually dimorphic projections to other peripheral targets \citep{adler2012neurochemical} and remain to be tested. Third, our survival interval was fixed at 6 days on the basis of pilot viral-progression studies; PRV labeling at different post-injection intervals would further disambiguate disynaptic from polysynaptic connectivity, as would dual-label tracing with a separate sensory tracer. Fourth, PRV152 reports connectivity but not function; whether the labeled neurons drive bone remodeling or bone pain in a behaviorally meaningful way requires combining this atlas with targeted activity perturbation, for example via PDE5A-defined or LepRb-defined neuronal subsets \citep{kim2020pde5a,flak2016cns}. Fifth, our injections were directed to the right femur only, and generalizing to other skeletal sites---vertebrae, skull, or ribs---is not directly supported by the present data.

Several implications for future work follow from this atlas. The most immediate is that it provides a prioritized list of central targets for causal perturbation experiments in bone biology, led by PAG, RMg, RPa, PVH, LH and MPOM. A second is that pharmacological strategies aimed at bone pain should consider that many of the same central neurons that gate nociception also deliver sympathetic drive to bone; centrally active analgesics that engage PAG or raphe circuitry may therefore have measurable downstream effects on bone metabolism, and vice versa. A third is that cell-type-specific reagents, for example Cre-driver lines targeting LepRb-, PDE5A-, or serotonergic raphe neurons, can now be used to decompose the atlas into functionally defined modules. Finally, extending the same tracing strategy to aged mice, to females, and to models of osteoporosis or cancer-induced bone pain would test the stability of this map across clinically relevant contexts. We believe this comprehensive atlas will stimulate further study of the neural regulation of bone metabolism and of bone pain.

\bibliographystyle{plainnat}
\bibliography{references}

\end{document}
